\section{Технический проект}

\subsection{Общая характеристика организации решения задачи}

Разрабатываемая система состоит из двух взаимосвязанных компонентов: библиотеки регулярных выражений на языке Си и десктопного поискового приложения на C++ с использованием фреймворка Qt.

Библиотека регулярных выражений реализована в виде единственного файла исходного кода на языке Си и заголовочного файла с описанием экспортируемого C API. Такой способ поставки обеспечивает максимальную переносимость: библиотека может быть включена в любой проект на языке Си или C++ без предварительной сборки в отдельный бинарный модуль. Заголовочный файл содержит объявления всех экспортируемых функций, обернутые в блок \verb`extern "C"`, что обеспечивает корректное связывание при компиляции в составе C++-проекта.

Поисковое приложение реализовано на C++ с использованием фреймворка Qt и включает в себя текстовый редактор с поддержкой поиска и замены, а также модуль рекурсивного поиска по файлам в директории. Приложение интегрирует библиотеку регулярных выражений напрямую через ее C API, подключая заголовочный файл и компилируя исходный файл библиотеки совместно с исходными файлами приложения.

\subsection{Архитектура системы}

Система разделена на два независимых компонента -- библиотеку регулярных выражений и поисковое приложение, -- каждый из которых включает в себя несколько модулей.

Библиотека регулярных выражений состоит из следующих модулей:
\begin{enumerate}
	\item Модуль парсера -- разбор строки шаблона регулярного выражения
	и построение внутреннего представления.
	\item Модуль сопоставления -- проверка входной строки на соответствие
	скомпилированному объекту регулярного выражения.
	\item Модуль замены -- выполнение замены совпадений в строке.
	\item Модуль управления памятью -- выделение и освобождение
	внутренних структур данных.
	\item Модуль автотестов -- набор автоматизированных тестов
	для верификации корректности работы библиотеки.
\end{enumerate}

Поисковое приложение состоит из следующих модулей:
\begin{enumerate}
	\item \verb`main` -- точка входа в приложение, инициализация Qt и запуск главного окна.
	\item \verb`TextEditor` -- модуль текстового редактора: открытие, отображение и редактирование файлов, подсветка совпадений, поиск и замена в тексте.
	\item \verb`SearchDialog` -- диалог поиска и замены в текстовом редакторе.
	\item \verb`SearchDir` -- модуль рекурсивного поиска по файлам в директории с отображением результатов и возможностью открытия найденного файла в редакторе.
\end{enumerate}


\subsection{Диаграмма компонентов}

На рисунке~\ref{diag:components} изображена диаграмма компонентов для проектируемой системы. Она включает модули:
\begin{itemize}
	\item модуль парсера;
	\item модуль сопоставления;
	\item модуль поиска и валидации;
	\item модуль замены;
	\item модуль управления памятью;
	\item модуль автотестов;
	\item главный модуль приложения (main);
	\item модуль текстового редактора (TextEditor);
	\item модуль диалога поиска и замены (SearchDialog);
	\item модуль рекурсивного поиска по директории (SearchDir).
\end{itemize}

\begin{figure}[H]
	\centering
	\includegraphics[width=0.95\linewidth]{Диаграмма компонентов.eps}
	\caption{Диаграмма компонентов проектируемой системы}
	\label{diag:components}      
\end{figure}

\subsection{Библиотека регулярных выражений}

\subsubsection{Внутренние структуры данных}

Скомпилированное регулярное выражение представляется в памяти в виде связного списка объектов типа \verb`struct re_chain` -- звеньев цепочки. Каждое звено соответствует одной альтерации регулярного выражения (разделяемых в шаблоне \verb`|`). Внутри звена хранится односвязный список объектов типа \verb`struct re_block` -- блоков, каждый из которых соответствует одному токену шаблона. Структура \verb`struct re_chain` содержит следующие поля:

\begin{itemize}
	\item \verb`next` -- указатель на следующее звено цепочки (следующую альтерацию) или \verb`NULL`;
	\item \verb`head`, \verb`tail` -- указатели на первый и последний блок в списке блоков данной альтернативы;
	\item \verb`bl_count` -- число блоков в списке;
	\item \verb`incl_start`, \verb`incl_end` -- флаги наличия якорей \verb`^` и \verb`$` в данной альтернативе;
	\item \verb`ch_data` -- для корневой цепочки: суммарное число групп захвата в регулярном выражении; для вложенных цепочек: порядковый номер соответствующей группы захвата;
	\item \verb`flags` -- модификаторы регулярного выражения, хранятся только в корневом звене цепочки; ненулевое значение является признаком того, что данное звено является корневым.
\end{itemize}

Структура \verb`struct re_block` содержит следующие поля:

\begin{itemize}
	\item \verb`type` -- тип блока, определяющий интерпретацию
	поля \verb`data`. Возможные типы блоков приведены
	в таблице~\ref{tab:block_types};
	\item \verb`data` -- объединение (\verb`union`) с данными блока в зависимости от его типа (см. таблицу ~\ref{tab:block_types});
	\item \verb`rep_min`, \verb`rep_max` -- минимальное и максимальное число повторений блока согласно квантификатору;
	\item \verb`lazy` -- флаг ленивого режима захвата;
	\item \verb`next` -- указатель на следующий блок в цепочке.
\end{itemize}

\begin{xltabular}{\textwidth}{|C{3cm}|T|C{5cm}|}
	\caption{Типы блоков регулярного выражения}
	\label{tab:block_types} \\
	\hline
	\thead{Тип} & \thead{Описание} & \thead{Данные блока} \\
	\hline
	\finishhead
	\verb`PLAIN` & Последовательность символов, воспринимаемых буквально;
	квантификаторы к блоку не применяются & Указатель на строку \verb`text` \\
	\hline
	\verb`ONE` & Одиночный символ, возможно с квантификатором &  Символ \verb`sym` \\
	\hline
	\verb`ANY` & Метасимвол «.» -- любой одиночный символ & Отсутствуют \\
	\hline
	\verb`CHARCLASS` & Символьный класс (\verb`\d`, \verb`\w`, \verb`\s` и др.) & Символ \verb`sym`, представляющий символьный класс \\
	\hline
	\verb`GROUP` & Набор символов в квадратных скобках \verb`[...]` & Указатель на строку \verb`text` с набором символов \\
	\hline
	\verb`REVGROUP` & Инвертированный набор символов \verb`[^...]` & Указатель на строку \verb`text` с набором символов \\
	\hline
	\verb`CHAIN` & Группа в круглых скобках & Указатель на вложенную цепочку альтераций \verb`subchain` \\
	\hline
	\verb`BACKREF` & Обратная ссылка на захваченную группу от \verb`\1` до \verb`\9` & Номер группы захвата \verb`cap_num` \\
	\hline
\end{xltabular}

Для хранения информации о захваченных подстроках в процессе сопоставления используется вспомогательная структура \verb`struct re_matches`, содержащая массивы начальных позиций (\verb`starts`) и длин (\verb`lengths`) захватов, их емкость (\verb`capacity`) и текущие флаги модификаторов (\verb`flags`).

\subsubsection{Внутреннее представление скомпилированного шаблона}

Результатом компиляции шаблона регулярного выражения является абстрактное синтаксическое дерево, реализованное в виде двух взаимосвязанных односвязных списков. Верхний уровень образуют звенья типа \verb`struct re_chain`, каждое из которых соответствует одной альтернативе шаблона (разделяемые оператором \verb`|`); звенья связаны через поле \verb`next`. Внутри каждого звена хранится список токенов конкатенации -- объектов типа \verb`struct re_block`, также связанных через поле \verb`next`. Блок типа \verb`CHAIN` (скобочная группа) рекурсивно содержит указатель на вложенную цепочку альтераций.

Корневое звено цепочки отличается от вложенных ненулевым значением поля \verb`flags`, которое хранит модификаторы регулярного выражения (\verb`RE_ROOT`, \verb`RE_M`, \verb`RE_S`, \verb`RE_I`, \verb`RE_IRU`). В поле \verb`ch_data` корневого звена хранится суммарное число групп захвата в шаблоне; в поле \verb`ch_data` вложенных звеньев -- порядковый номер соответствующей группы захвата.

На рисунке~\ref{diag:ast_diagram} приведена диаграмма внутреннего представления для шаблона \verb`(a+|^bc)\d{2}`. Корневая цепочка содержит два блока: блок типа \verb`CHAIN` (группа захвата~№1) с квантификатором \verb`{1,1}` и блок типа \verb`CHARCLASS` (символьный класс \verb`\d`) с квантификатором \verb`{2,2}`. Блок \verb`CHAIN` ссылается на вложенную цепочку из двух
альтернатив: первая содержит блок типа \verb`ONE` (символ \verb`a` с квантификатором \verb`{1,∞}`), вторая -- с флагом захвата с начала строки и блоком типа \verb`PLAIN` (буквальная строка \verb`"bc"`).

\begin{figure}[H]
	\centering
	\includegraphics[width=0.95\linewidth]{Диаграмма АСД.eps}
	\caption{Диаграмма абстрактного синтаксического дерева регулярного выражения}
	\label{diag:ast_diagram}      
\end{figure}

\subsubsection{Формальная грамматика шаблона регулярного выражения}

Синтаксис шаблонов регулярных выражений, поддерживаемых библиотекой, описывается формальной грамматикой в форме Бэкуса-Наура (BNF). Грамматика отражает приоритет операций: квантификаторы связывают сильнее конкатенации, конкатенация -- сильнее альтерации.

\begin{lstlisting}[language={},basicstyle=\linespread{1.1}\ttfamily,frame=single,caption={Грамматика шаблона регулярного выражения в форме Бэкуса--Наура}]
	<шаблон>          ::= <альтерация>
	<альтерация>      ::= <конкатенация>
	| <альтерация>  "|" <конкатенация>
	<конкатенация>    ::= <токен>
	| <конкатенация> <токен>
	<токен>           ::= <атом>
	| <атом> <квантификатор>
	| <атом> <квантификатор> "?"
	<квантификатор>   ::= "*" | "+" | "?"
	| "{" <число> "}"
	| "{" <число> "," <число> "}"
	<атом>            ::= <символ>
	| "\" <метасимвол>
	| "\" <символьный-класс>
	| "\" <цифра-от-1-до-9>
	| "."
	| "^"
	| "$"
	| <набор-символов>
	| <группа>
	<группа>          ::= "(" <шаблон> ")"
	| "(?:" <шаблон> ")"
	| "(?flags)"
	<набор-символов>  ::= "[" <элементы-набора> "]"
	| "[^" <элементы-набора> "]"
	<элементы-набора> ::= <элемент-набора>
	| <элемент-набора> <элементы-набора>
	<элемент-набора>  ::= <символ-набора>
	| "\" <символьный-класс>
	| "\" <символ-набора>
	| <символ-набора> "-" <символ-набора>
	<символьный-класс>::= "d" | "D" | "w" | "W" | "s" | "S" | "a" | "A" | "p" | "P" | "n" | "t"
	<модификатор>     ::= "m" | "s" | "i" | "I"
	<символ>          ::= любой символ ASCII или байтовая последовательность
	UTF-8, кроме метасимволов
	<метасимвол>      ::= "." | "^" | "$" | "*" | "+" | "?" | "[" | "]" | "{" | "}" | "|" | "(" | ")" | "\"
	<символ-набора>   ::= любой символ, кроме "]" и "\"
	<цифра-от-1-до-9> ::= "1" | "2" | "3" | "4" | "5" | "6" | "7" | "8" | "9"
	<число>           ::= <цифра> | <число> <цифра>
	<цифра>           ::= "0" | "1" | "2" | "3" | "4" | "5" | "6" | "7" | "8" | "9"
\end{lstlisting}

Несколько пояснений к грамматике:
\begin{itemize}
	\item форма \verb`(?flags)` задает модификаторы регулярного выражения
	(\verb`m`, \verb`s`, \verb`i`, \verb`I`) и не производит захвата;
	обратная ссылка \verb`\n` (где \textit{n} -- цифра от 1 до 9)
	ссылается на ранее захваченную группу и допустима только в позиции атома;
	\item в строке замены допускаются ссылки вида \verb`$n` или \verb`\n`
	(где \textit{n} от 0 до 9), которые не являются частью шаблона
	и в грамматике не описываются;
	\item диапазон \verb`{n,m}` требует $n \leq m$; оба числа не должны
	превышать \verb`INT_MAX`.
\end{itemize}

\subsubsection{Алгоритм компиляции шаблона}

Компиляция шаблона выполняется рекурсивным нисходящим парсером, реализованным в функции \verb`_parse_chain`. Результатом компиляции является абстрактное синтаксическое дерево (АСД) -- иерархическая структура, в которой альтерации представлены звеньями списка \verb`re_chai`n, а токены каждой альтернативы -- блоками \verb`re_block`, образующими вложенное дерево через блоки типа \verb`CHAIN`. АСД формируется непосредственно в процессе однопроходного разбора строки шаблона без построения промежуточных представлений, что сокращает число выделений памяти и упрощает реализацию.

При встрече с открывающей круглой скобкой парсер рекурсивно вызывает себя для разбора вложенной группы, результат которого записывается в блок типа \verb`CHAIN`. При встрече с символом альтерации «\verb`|`» создается новое звено цепочки, и разбор продолжается в нем. Квантификаторы (\verb`*`, \verb`+`, \verb`?`, \verb`{n,m}`) обрабатываются вспомогательной функцией \verb`process_quantif` непосредственно после распознавания предшествующего токена.

Проверка корректности шаблона выполняется в процессе разбора без отдельного прохода. При обнаружении синтаксической ошибки парсер немедленно прерывает работу через механизм нелокального перехода \verb`setjmp`/\verb`longjmp`, возвращая управление в функцию \verb`compile_re2` с кодом ошибки и позицией символа в шаблоне, при разборе которого ошибка была обнаружена.

Если в шаблоне заданы флаги регистронезависимого поиска (\verb`i` или \verb`I`), после завершения разбора над всем внутренним представлением выполняется однократный проход \verb`re_tolower_inplace`, приводящий все литеральные данные в блоках к нижнему регистру. В процессе сопоставления аналогичное приведение выполняется над входной строкой, что позволяет реализовать регистронезависимый поиск без изменения алгоритма сопоставления.

Точкой входа при компиляции является экспортируемая функция \verb`compile_re2`, вызывающая \verb`_parse_chain` с начальными параметрами. Упрощенная обертка \verb`compile_re` вызывает \verb`compile_re2` с \verb`NULL`-указателями вместо параметров диагностики ошибок.

На рисунке ~\ref{diag:activity_diag} изображена диаграмма активностей алгоритма компиляции шаблона регулярного выражения.
\begin{figure}[H]
	\centering
    \includegraphics[width=0.95\linewidth]{Диаграмма активностей.eps}
    \caption{Диаграмма активностей алгоритма компиляции}
    \label{diag:activity_diag}      
\end{figure}

\subsubsection{Алгоритм сопоставления с образцом}

Для проверки соответствия входной строки скомпилированному регулярному выражению применяется алгоритм поиска с возвратами (метод перебора с откатом). Алгоритм реализован в функции \verb`_match_chain` и работает следующим образом.

Для каждой стартовой позиции во входной строке алгоритм последовательно обходит блоки текущего звена цепочки. Для каждого блока он жадно (или лениво, если задан флаг \verb`lazy`) пытается захватить максимально возможное число повторений, соответствующих квантификатору. Если следующий блок не сопоставляется, алгоритм возвращается к предыдущему блоку и уменьшает число захваченных им повторений на единицу, после чего предпринимает новую попытку. Таким образом осуществляется полный перебор допустимых комбинаций с откатами.

Для блоков типа \verb`CHAIN` (вложенных групп) сопоставление выполняется рекурсивно через функцию \verb`_match_group`, которая перебирает все альтернативы вложенной цепочки и выбирает ту из них, начало захвата которой расположено как можно раньше во входной строке, -- в соответствии с поведением регулярных выражений в других языках.

Позиции и длины захваченных групп в процессе сопоставления накапливаются в структуре \verb`struct re_matches` и используются как для возврата результата вызывающей стороне, так и для разрешения обратных ссылок (\verb`BACKREF`) в ходе того же сопоставления.

Подготовка к сопоставлению (нормализация входной строки при регистронезависимом поиске, обработка начальной позиции поиска) выполняется оберткой \verb`_match_wrapper`, которую вызывают все экспортируемые функции сопоставления.

\subsubsection{Коды ошибок компиляции}

При обнаружении синтаксической ошибки в шаблоне регулярного выражения функция \verb`compile_re2` возвращает \verb`NULL` и записывает по указателю \verb`errcode` один из кодов ошибки, перечисленных в таблице~\ref{tab:errorcodes}. По указателю \verb`errpos` записывается индекс символа в строке шаблона, при разборе которого ошибка была обнаружена. Человекочитаемое описание ошибки на английском языке можно получить вызовом отладочной функции \verb`get_error_msg(errcode)`.

\begin{xltabular}{\textwidth}{|C{2cm}|C{5cm}|T|}
	\caption{Коды ошибок компиляции регулярного выражения}
	\label{tab:errorcodes} \\
	\hline
	\thead{Код} & \thead{Константа} & \thead{Описание} \\
	\hline
	\finishhead
	$-100$ & \verb`BAD_CARET` & Неэкранированный символ \verb`^` не в начале регулярного выражения или набора символов \\
	\hline
	$-99$ & \verb`BAD_DOLLAR` & Неэкранированный символ \verb`$` не в конце регулярного выражения \\
	\hline
	$-98$ & \verb`PREV_NOT_QUANTIF` & Предыдущий токен не квантифицируем \\
	\hline
	$-97$ & \verb`EMPTY_SQUARES` & Пустой набор символов в квадратных скобках \\
	\hline
	$-96$ & \verb`MISSING_LSQUARE` & Отсутствует открывающая квадратная скобка «\verb`[`» \\
	\hline
	$-95$ & \verb`MISSING_RSQUARE` & Отсутствует закрывающая квадратная скобка «\verb`]`» \\
	\hline
	$-94$ & \verb`RANGE_WITH_CHARCLASS` & Диапазон с символьным классом недопустим (например, \verb`[\w-z]`) \\
	\hline
	$-93$ & \verb`BAD_ESC_CHAR` & Отсутствует символ после экранирующего символа \\
	\hline
	$-92$ & \verb`BAD_INSIDE_CURLY` & Недопустимые символы внутри квантификатора в фигурных скобках \\
	\hline
	$-91$ & \verb`BAD_CURLY_ORDER` & Неверный порядок чисел в квантификаторе в фигурных скобках \\
	\hline
	$-90$ & \verb`HUGE_CURLY_RANGE` & Диапазон квантификатора слишком велик \\
	\hline
	$-89$ & \verb`MISSING_LPAREN` & Отсутствует открывающая круглая скобка «\verb`(`» \\
	\hline
	$-88$ & \verb`MISSING_RPAREN` & Отсутствует закрывающая круглая скобка «\verb`)`» \\
	\hline
	$-87$ & \verb`BAD_SQUARE_ORDER` & Неверный порядок символов в диапазоне набора символов \\
	\hline
	$-86$ & \verb`BAD_CHARCLASS` & Некорректный (несуществующий) символьный класс \\
	\hline
	$-85$ & \verb`INVALID_FLAG` & Неверный модификатор регулярного выражения \\
	\hline
\end{xltabular}


\subsubsection{Алгоритмы экспортируемых функций}

Все экспортируемые функции библиотеки построены по единой схеме: публичная функция является тонкой обёрткой над внутренней реализацией, принимающей как строку шаблона, так и скомпилированный объект. Если передана строка шаблона, внутренняя функция компилирует её через \verb`compile_re2`, выполняет операцию и немедленно освобождает объект через \verb`free_re`. Если передан готовый объект, этапы компиляции и освобождения пропускаются. Такая схема позволяет избежать дублирования логики между \verb`pattern`- и \verb`obj_`-версиями каждой функции.

\subsubsubsection{Семейство \texttt{matchpos}}

Внутренняя функция \verb`_matchpos` принимает либо строку шаблона, либо объект регулярного выражения. При получении строки шаблона она компилирует его через \verb`compile_re2`; при ошибке компиляции записывает позицию ошибочного символа по указателю \verb`capture_data` и возвращает код ошибки. При успешной компиляции вызывает \verb`_match_wrapper`, передавая значение \verb`*capture_data` как начальную позицию поиска, если оно положительно. Возвращает начальную позицию совпадения и записывает его длину по указателю \verb`capture_data`. Если совпадений не найдено, возвращает $-1$.

\subsubsubsection{Семейство \texttt{is\_match}}

Внутренняя функция \verb`_is_match` аналогична \verb`_matchpos`, но передаёт в \verb`_match_wrapper` \verb`NULL` вместо указателя на длину захвата, поскольку длина результата не требуется. Возвращаемое значение нормализуется: $-1$ от \verb`_match_wrapper` (совпадений нет) преобразуется в $0$, любое неотрицательное значение -- в $1$.

\subsubsubsection{Семейство \texttt{match}}

Внутренняя функция \verb`_match` выделяет структуру \verb`re_matches` с массивами начальных позиций и длин групп захвата ёмкостью, равной числу групп захвата в скомпилированном объекте. После вызова \verb`_match_wrapper` формирует результирующий массив: первым элементом записывается \verb`malloc`-строка с полным совпадением, затем последовательно -- строки для каждой группы захвата с непустым результатом. Массив завершается указателем на \verb`NULL`. При отсутствии совпадений массив содержит только \verb`NULL`. Временная структура \verb`re_matches` освобождается до возврата.

\subsubsubsection{Семейство \texttt{matches}}

Внутренняя функция \verb`_matches` выполняет итеративный поиск всех непересекающихся совпадений. На каждой итерации вызывается \verb`_match_wrapper` с текущей позицией поиска \verb`pos`. При успешном захвате подстрока копируется в \verb`malloc`-строку и добавляется в результирующий массив; позиция \verb`pos` сдвигается за конец найденного совпадения. При нулевом захвате позиция сдвигается на один символ, чтобы исключить зависание на месте. Цикл завершается при первом неудачном поиске или при достижении конца строки.

Результирующий массив строк реализован с динамическим наращиванием ёмкости: начальная ёмкость равна $10$ элементам; при её исчерпании массив перевыделяется через \verb`re_realloc` с удвоением ёмкости. Такая стратегия обеспечивает амортизированную стоимость добавления элемента $O(1)$ и позволяет не знать заранее число совпадений. Массив завершается указателем \verb`NULL`.

\subsubsubsection{Семейства \texttt{re\_nreplace} и \texttt{re\_dupreplace}}

Оба семейства реализованы через единую внутреннюю функцию \verb`_re_replace`, поведение которой определяется значением указателя \verb`*dest_ptr`: если он равен \verb`NULL`, активируется режим динамического выделения памяти (\verb`re_dupreplace`), иначе результат записывается в буфер фиксированного размера (\verb`re_nreplace`).

Алгоритм замены работает в цикле: на каждой итерации вызывается \verb`_match_wrapper` с текущей позицией \verb`curr_char`. При нахождении совпадения в результирующий буфер сначала копируется префикс -- часть исходной строки от \verb`curr_char` до начала совпадения, -- затем обрабатывается строка замены \verb`repl`. В строке замены распознаются обратные ссылки вида \verb`$n` и \verb`\n` (где $n$ от 0 до 9): \verb`$0` заменяется полным совпадением, \verb`$1`--\verb`$9` -- содержимым соответствующей группы захвата, если она была захвачена. Позиции и длины групп захвата хранятся в локальной структуре \verb`re_matches` с массивами \verb`starts` и \verb`lengths` ёмкостью 10 элементов, расположенными на стеке. При нулевом захвате в буфер дополнительно записывается символ исходной строки на позиции совпадения, после чего позиция сдвигается на единицу (поведение при нулевом захвате не стандартизировано и различается между реализациями; в данной библиотеке принята модель PCRE/JavaScript как наиболее распространённая среди промышленных библиотек). После завершения цикла в буфер дописывается остаток исходной строки, начиная с \verb`curr_char`.

В режиме фиксированного буфера копирование на каждом шаге ограничивается числом оставшихся байтов; при исчерпании буфера функция возвращает $-1$. В режиме динамического выделения начальный размер буфера равен длине исходной строки плюс один байт под нуль-терминатор. При нехватке места буфер перевыделяется через \verb`re_realloc` с запасом, равным длине исходной строки, -- это снижает число перевыделений при многократных заменах на строки, длиннее исходных совпадений. Функция возвращает число выполненных замен, а указатель на выделенный буфер записывается по \verb`dest_ptr`.

\subsubsection{Управление памятью и правила владения}

Библиотека придерживается явных и однозначных правил владения памятью, которые должны соблюдаться вызывающей стороной.

Объекты регулярных выражений, созданные функциями \verb`compile_re` и \verb`compile_re2`, должны быть освобождены вызовом \verb`free_re` после окончания использования. Библиотека не хранит глобальных указателей на такие объекты и не освобождает их самостоятельно, кроме случаев внутреннего использования в \verb`pattern`-версиях функций.

Массивы строк, возвращаемые функциями \verb`match` и \verb`matches`, выделены через \verb`malloc` в два уровня: сам массив указателей и каждая строка в нём. Вызывающая сторона обязана освободить каждую строку через \verb`free`, а затем освободить массив указателей. Строки, возвращаемые \verb`re_dupreplace` и \verb`obj_re_dupreplace`, также выделены через \verb`malloc` и должны быть освобождены вызывающей стороной.

Буфер \verb`dest` в функциях \verb`re_nreplace` и \verb`obj_re_nreplace` предоставляется вызывающей стороной и библиотекой не освобождается.

Все внутренние выделения памяти выполняются через обёртки \verb`re_malloc` и \verb`re_realloc`, завершающие процесс с сообщением об ошибке при исчерпании памяти, что исключает необходимость проверки возвращаемых значений внутри библиотеки. Скомпилированный объект \verb`struct re_chain` в процессе сопоставления не модифицируется -- все промежуточные данные о захватах хранятся в структуре \verb`struct re_matches`, создаваемой на стеке при каждом вызове. Глобального изменяемого состояния библиотека не использует. Поэтому параллельное использование одного скомпилированного объекта из нескольких потоков безопасно без какой-либо синхронизации. Независимые объекты, созданные в разных потоках, также могут использоваться параллельно без ограничений.

\subsection{Поисковое приложение}

\subsubsection{Общая архитектура приложения}

Поисковое приложение построено по архитектурному шаблону, принятому во фреймворке Qt для приложений с графическим интерфейсом: каждому окну или диалогу соответствует пара классов -- класс окна и автоматически генерируемый из \verb`.ui`-файла класс \verb`Ui::*`, инкапсулирующий виджеты интерфейса. Бизнес-логика сосредоточена в классах окон; прямое обращение к виджетам из внешних модулей не допускается.

Приложение работает в двух режимах: режиме текстового редактора с поиском и заменой в открытом файле и режиме поиска по директории. Переключение между режимами осуществляется открытием соответствующего диалогового окна из главного меню редактора.

Интеграция библиотеки регулярных выражений выполнена на уровне исходного кода: файлы \verb`regex-lib.c` и \verb`regex-lib.h` включены непосредственно в проект Qt и компилируются вместе с исходными файлами приложения. Заголовочный файл подключается через блок \verb`extern "C"` во всех модулях, использующих API библиотеки.

Приложение состоит из пяти программных модулей, описанных в таблице~\ref{tab:app_modules}.

\begin{xltabular}{\textwidth}{|C{3.5cm}|C{3.5cm}|T|}
	\caption{Модули поискового приложения}
	\label{tab:app_modules} \\
	\hline
	\thead{Модуль} & \thead{Класс Qt} & \thead{Назначение} \\
	\hline
	\finishhead
	\verb`main` & -- & Точка входа; создаёт \verb`QApplication` и главное окно \verb`TextEditor` \\
	\hline
	\verb`texteditor` & \verb`QMainWindow` & Главное окно: текстовый редактор, меню, строка состояния, масштабирование \\
	\hline
	\verb`searchdialog` & \verb`QDialog` & Немодальный диалог поиска и замены в открытом файле \\
	\hline
	\verb`searchdir` & \verb`QDialog` & Немодальный диалог рекурсивного поиска по файлам директории \\
	\hline
	\verb`searchworker` & \verb`QObject`, \verb`QRunnable` & Фоновые классы поиска: \verb`SearchController` управляет обходом директории, \verb`FileSearchWorker` выполняет поиск в одном файле \\
\end{xltabular}

\subsubsection{Модуль \texttt{TextEditor}}

Класс \verb`TextEditor` наследует \verb`QMainWindow` и является главным окном приложения. Центральным виджетом служит \verb`QPlainTextEdit` -- \verb`text_edit_main`, в котором отображается и редактируется содержимое открытого файла. Строка состояния \verb`QStatusBar` отображает текущую позицию курсора (строка и столбец), тип переноса строк и масштаб отображения текста.

Модуль реализует следующие функции.

\textit{Открытие файла.} Слот \verb`on_action_open_triggered` открывает диалог выбора файла \verb`QFileDialog`. Файл читается как байтовый массив; двоичные файлы отклоняются с диалогом ошибки. Кодировка определяется автоматически: при обнаружении файла с кодировкой CP1251 происходит его конвертация в UTF-8 для дальнейшей работы. Тип переноса строк (\verb`LF`, \verb`CR`, \verb`CRLF`) определяется функцией \verb`detectLineEnding` и сохраняется для последующей записи.

\textit{Сохранение файла.} Слот \verb`on_action_save_triggered` сохраняет текущее содержимое редактора в исходный файл с восстановлением оригинального типа переноса строк. Слот \verb`on_action_save_as_triggered` открывает диалог сохранения \verb`QFileDialog` и позволяет выбрать новое имя или путь.

\textit{Масштабирование.} Слоты \verb`on_action_zoom_in_triggered`, \verb`on_action_zoom_out_triggered` и \verb`on_action_zoom_100_triggered` изменяют размер шрифта редактора по списку стандартных значений \verb`ZOOM_SIZES`. Масштабирование также доступно через \verb`Ctrl+Колёсико мыши` посредством фильтра событий \verb`eventFilter`. Текущий масштаб в процентах отображается в строке состояния.

\textit{Подсветка совпадений.} Для подсветки найденных совпадений применяется механизм дополнительных выделений \verb`QPlainTextEdit::setExtraSelections`, управляемый классом \verb`SearchDialog`. Список выделений сбрасывается при изменении текста через соединение сигнала \verb`textChanged` со слотом \verb`SearchDialog::clearHighlights`.

\textit{Отображение позиции курсора.} Сигнал \verb`cursorPositionChanged` редактора соединён со слотом \verb`onCursorPositionChanged`, который вычисляет текущие номер строки и столбец через \verb`QTextCursor::blockNumber` и \verb`QTextCursor::positionInBlock` и отображает их в постоянном виджете \verb`status_cursor` строки состояния.

\textit{Взаимодействие с диалогом поиска.} Объект \verb`SearchDialog` создаётся при инициализации \verb`TextEditor` и существует на протяжении всего времени жизни главного окна. Диалог является немодальным. Слот \verb`on_action_find_triggered` показывает диалог; при первом показе диалог позиционируется справа от главного окна, если позволяет место на экране, иначе -- по центру экрана.

Поля класса \verb`TextEditor` приведены в таблице~\ref{tab:texteditor_fields}.

\begin{xltabular}{\textwidth}{|C{4.5cm}|C{3cm}|T|}
	\caption{Поля класса \texttt{TextEditor}}
	\label{tab:texteditor_fields} \\
	\hline
	\thead{Поле} & \thead{Тип} & \thead{Назначение} \\
	\hline
	\finishhead
	\verb`ui` & \verb`Ui::TextEditor*` & Указатель на виджеты, сгенерированные из \verb`.ui`-файла \\
	\hline
	\verb`search_dialog` & \verb`SearchDialog*` & Немодальный диалог поиска и замены \\
	\hline
	\verb`search_dir_dialog` & \verb`SearchDir*` & Немодальный диалог поиска по директории; \verb`nullptr` если не открыт \\
	\hline
	\verb`current_file_path` & \verb`QString` & Путь к текущему открытому файлу; пуст, если файл не открыт \\
	\hline
	\verb`current_line_ending` & \verb`LineEnding` & Тип переноса строк открытого файла (\verb`LE_LF`, \verb`LE_CR`, \verb`LE_CRLF`) \\
	\hline
	\verb`current_zoom_percent` & \verb`int` & Текущий масштаб отображения в процентах \\
	\hline
	\verb`status_cursor` & \verb`QLabel*` & Постоянный виджет строки состояния: позиция курсора \\
	\hline
	\verb`status_line_ending` & \verb`QLabel*` & Постоянный виджет строки состояния: тип переноса строк \\
	\hline
	\verb`status_zoom` & \verb`QLabel*` & Постоянный виджет строки состояния: масштаб \\
\end{xltabular}

\subsubsection{Модуль \texttt{SearchDialog}}

Класс \verb`SearchDialog` наследует \verb`QDialog` и реализует немодальный диалог поиска и замены. Диалог принимает указатель на \verb`QPlainTextEdit` главного окна и указатель на \verb`QStatusBar` и напрямую работает с содержимым и курсором редактора.

Диалог поддерживает два режима поиска, выбираемых переключателями (\verb`QRadioButton`): обычный текстовый поиск средствами \verb`QTextDocument::find` и поиск по регулярному выражению с использованием API библиотеки \verb`regex-lib`. В режиме регулярных выражений дополнительно доступны опции многострочного режима (\verb`(?m)`) и режима "<\textquotedbl .\textquotedbl \ -- новая строка"> (\verb`(?s)`); опция поиска целых слов в этом режиме недоступна и блокируется автоматически.

Область поиска может быть ограничена текущим выделением через флажок «Только в выделении». Поддерживается циклический поиск: при достижении конца текста поиск продолжается с начала с уведомлением в строке статуса.

Панель замены скрыта по умолчанию и разворачивается кнопкой-переключателем. В режиме поиска вместо кнопки «Заменить всё» отображается кнопка «Подсветить» -- она подсвечивает все совпадения в области поиска через \verb`setExtraSelections`.

Модуль реализует следующие операции.

\textit{Поиск следующего совпадения.} Метод \verb`findNext_regex` определяет область поиска через \verb`getSearchScope`, вычисляет начальную позицию в байтах UTF-8 от текущего положения курсора и вызывает \verb`matchpos`. Функция \verb`selectInEditor` переводит байтовое смещение в позицию символа Qt и выделяет найденный фрагмент. При нулевом захвате выделяется ближайший символ с уведомлением в строке статуса.

\textit{Подсчёт и подсветка совпадений.} Метод \verb`collectMatches` итеративно вызывает \verb`matchpos` и заполняет список \verb`found_matches`. Метод \verb`highlightAll` формирует список \verb`QTextEdit::ExtraSelection` по этому списку и передаёт его в \verb`setExtraSelections`.

\textit{Замена одного совпадения.} Слот \verb`on_btn_replace_clicked` в режиме регулярных выражений вызывает \verb`re_dupreplace` с \verb`limit=1`, применяя замену к фрагменту текста начиная с позиции текущего выделения, после чего выполняет поиск следующего совпадения.

\textit{Замена всех совпадений.} Слот \verb`on_btn_replace_all_highlight_clicked` в режиме замены вызывает \verb`re_dupreplace` с \verb`limit=0` для области поиска и заменяет соответствующий фрагмент содержимого редактора результатом. При ограниченной области поиска результат склеивается с частями текста за её пределами.

\textit{Валидация шаблона.} Метод \verb`validatePattern` компилирует шаблон через \verb`compile_re2` и при обнаружении ошибки отображает HTML-форматированное сообщение с подчёркиванием ошибочного символа. Смещение позиции ошибки корректируется на длину автоматически добавленного префикса модификаторов.

Поля класса \verb`SearchDialog` приведены в таблице~\ref{tab:searchdialog_fields}.

\begin{xltabular}{\textwidth}{|C{4.5cm}|C{3.5cm}|T|}
	\caption{Поля класса \texttt{SearchDialog}}
	\label{tab:searchdialog_fields} \\
	\hline
	\thead{Поле} & \thead{Тип} & \thead{Назначение} \\
	\hline
	\finishhead
	\verb`ui` & \verb`Ui::SearchDialog*` & Указатель на виджеты, сгенерированные из \verb`.ui`-файла \\
	\hline
	\verb`text_editor` & \verb`QPlainTextEdit*` & Указатель на редактор главного окна \\
	\hline
	\verb`main_status_bar` & \verb`QStatusBar*` & Строка состояния главного окна \\
	\hline
	\verb`replace_visible` & \verb`bool` & Флаг видимости панели замены \\
	\hline
	\verb`found_matches` & \verb`QList<Match>` & Список найденных совпадений (байтовые позиции и длины) для подсветки \\
	\hline
	\verb`current_match_idx` & \verb`int` & Индекс текущего активного совпадения в \verb`found_matches`; $-1$ если нет \\
\end{xltabular}

\subsubsection{Модуль \texttt{SearchDir}}

Класс \verb`SearchDir` наследует \verb`QDialog` и реализует немодальный диалог рекурсивного поиска по регулярному выражению или обычному тексту во всех текстовых файлах выбранной директории. Диалог открывается из главного меню редактора после выбора директории.

Интерфейс диалога включает поле ввода пути к директории с кнопкой обзора, поле ввода шаблона, комбинированный список масок расширений файлов с предустановленными вариантами («Текстовые файлы», «Исходный код», «Все файлы»), флажки режимов поиска, кнопки запуска и остановки, прогресс-бар и дерево результатов \verb`QTreeWidget`.

Дерево результатов организовано двухуровневой иерархией: узлы первого уровня соответствуют файлам с найденными совпадениями и содержат относительный путь к файлу и число совпадений в нём; узлы второго уровня соответствуют отдельным совпадениям и содержат номер строки, номер столбца и фрагмент текста вокруг совпадения. При двойном щелчке по любому узлу соответствующий файл открывается в редакторе, а курсор переводится на строку найденного совпадения.

Поиск выполняется в фоновом потоке через класс \verb`SearchController`, перемещаемый в \verb`QThread`. \verb`SearchController` обходит директорию через \verb`QDirIterator` и создаёт объекты \verb`FileSearchWorker` -- по одному на каждый файл. Воркеры помещаются в локальный \verb`QThreadPool`; при включённом многопоточном режиме пул использует все доступные ядра процессора, при выключенном -- одно. Завершение поиска отслеживается через атомарный счётчик активных воркеров \verb`active_workers`.

Для автоматического исключения двоичных файлов из поиска применяется эвристический метод: из каждого файла считываются первые 8~КБ содержимого, и при обнаружении нулевого байта файл классифицируется как двоичный и пропускается. Данный подход основан на том, что текстовые файлы в кодировках ASCII, UTF-8 и CP1251 не содержат нулевых байтов, тогда как подавляющее большинство двоичных форматов содержит их в самом начале файла.

Остановка поиска реализована через атомарный флаг \verb`stop_flag` типа \verb`QAtomicInt`, проверяемый в каждой итерации воркера. Кнопка «Найти» блокируется на время завершения потока после запроса остановки и разблокируется по сигналу \verb`searchFinished`.

Поля класса \verb`SearchDir` приведены в таблице~\ref{tab:searchdir_fields}.

\begin{xltabular}{\textwidth}{|C{4.5cm}|C{3.5cm}|T|}
	\caption{Поля класса \texttt{SearchDir}}
	\label{tab:searchdir_fields} \\
	\hline
	\thead{Поле} & \thead{Тип} & \thead{Назначение} \\
	\hline
	\finishhead
	\verb`ui` & \verb`Ui::SearchDir*` & Указатель на виджеты, сгенерированные из \verb`.ui`-файла \\
	\hline
	\verb`text_editor` & \verb`TextEditor*` & Указатель на главное окно для открытия найденных файлов \\
	\hline
	\verb`re_obj` & \verb`struct re_chain*` & Скомпилированный объект регулярного выражения; передаётся воркерам и освобождается по завершении поиска \\
	\hline
	\verb`worker` & \verb`SearchController*` & Объект контроллера фонового поиска \\
	\hline
	\verb`thread` & \verb`QThread*` & Фоновый поток, в котором выполняется \verb`SearchController` \\
	\hline
	\verb`search_running` & \verb`bool` & Флаг активного поиска; управляет состоянием кнопки запуска/остановки \\
	\hline
	\verb`total_matches` & \verb`int` & Суммарное число найденных совпадений по всем файлам \\
\end{xltabular}

\subsubsection{Интеграция библиотеки регулярных выражений}

В модулях \verb`SearchDialog` и \verb`SearchDir` взаимодействие с библиотекой регулярных выражений организовано по единой схеме. При запуске поиска шаблон компилируется однократно через \verb`compile_re2`; полученный объект \verb`struct re_chain*` передается во все последующие вызовы \verb`obj_matchpos` и \verb`obj_re_dupreplace`. По завершении работы с объектом вызывается \verb`free_re`. Такой подход исключает повторную компиляцию одного и того же шаблона при итеративном поиске и особенно важен в модуле \verb`SearchDir`, где один шаблон применяется последовательно к большому числу файлов.

Передача текста между Qt и библиотекой выполняется через промежуточное UTF-8-представление: \verb`QString` преобразуется в \verb`QByteArray` методом \verb`toUtf8`, после чего указатель \verb`QByteArray::data` передается в функции библиотеки. Результирующие позиции совпадений, возвращаемые в байтовых смещениях UTF-8, преобразуются обратно в позиции символов Qt вспомогательной функцией \verb`selectInEditor` через \verb`QString::fromUtf8`.